% chapters/09/section-03.tex
\section{Menerapkan Mesin Keadaan di Dalam Kode}
\label{sec:menerapkan-fsm}

Diagram pada bagian sebelumnya diterjemahkan menjadi kode. Susunan yang dipakai
pada bab ini memisahkan tiga pekerjaan yang berbeda, dan pemisahan itulah yang
menjadi inti bagian ini.

\subsection{Tiga Pekerjaan yang Dipisahkan}

\begin{enumerate}
  \item \textbf{Menentukan keadaan berikutnya.} Dikerjakan oleh fungsi
        \code{decide}, yang membaca keadaan sekarang, jarak terhadap pemain,
        dan ketiga pewaktu, lalu mengembalikan nama keadaan berikutnya. Fungsi
        ini tidak mengubah apa pun.
  \item \textbf{Menyiapkan keadaan yang baru dimasuki.} Dikerjakan oleh fungsi
        \code{enter_state}, yang mengembalikan penghitung keadaan sebelumnya ke
        nol dan mencatat perpindahan itu pada catatan keadaan.
  \item \textbf{Menjalankan perilaku keadaan yang aktif.} Dikerjakan oleh fungsi
        \code{act}, yang mengubah kecepatan NPC menurut keadaan yang sedang
        aktif.
\end{enumerate}

Fungsi \code{decide} tidak mengubah apa pun; ia hanya membaca
dan mengembalikan jawaban. Sifat itu disengaja, dan manfaatnya akan terasa pada
Bagian~\ref{sec:menguji-keadaan}: keputusan dapat diperiksa dengan memanggil
fungsi itu berulang kali tanpa menjalankan permainannya.

\subsection{Syarat Perpindahan Ditulis pada Satu Tempat}

Seluruh syarat pada Tabel~\ref{tab:tujuh-perpindahan} dikumpulkan di dalam
fungsi \code{decide}, disusun menurut keadaan yang sedang aktif.

\begin{kodekeluaran}
def decide(self, delta_seconds):
    self.attack_timer = max(0.0, self.attack_timer - delta_seconds)

    if self.state == IDLE_STATE:
        self.idle_timer += delta_seconds
        if self.distance <= self.radius:
            return CHASE_STATE
        if self.idle_timer >= IDLE_SECONDS:
            return PATROL_STATE
        return IDLE_STATE
    # ... keadaan berikutnya ...
\end{kodekeluaran}

Tiga hal perlu diperhatikan pada susunan itu.

\begin{enumerate}
  \item \textbf{Setiap keadaan diperiksa secara terpisah}, sehingga syarat
        perpindahannya mudah dicari. Menambahkan keadaan baru berarti
        menambahkan satu blok, bukan menyunting blok yang sudah ada.
  \item \textbf{Urutan pemeriksaan menentukan hasil} ketika dua syarat benar
        bersamaan. Pada keadaan berdiam, syarat ``pemain terlihat'' diperiksa
        lebih dahulu daripada syarat ``sudah cukup lama berdiam''. Pilihan itu
        disengaja: NPC yang melihat pemain harus segera mengejar, walaupun ia
        baru saja mulai berdiam.
  \item \textbf{Ketiga pewaktu dijalankan di tempat yang berbeda menurut
        jenisnya.} Pewaktu jeda serangan dijalankan pada awal fungsi, sebab ia
        berjalan tanpa memandang keadaan. Pewaktu berdiam dijalankan pada blok
        keadaan berdiam, dan pewaktu tenggang pada blok keadaan mengejar.
\end{enumerate}

\subsection{Pekerjaan Penyiapan ketika Masuk Keadaan}

Perpindahan keadaan bukan sekadar mengganti nama. Setiap keadaan memiliki
pekerjaan penyiapan sendiri, dan pekerjaan itu dikerjakan sekali saja pada saat
keadaan dimasuki.

\begin{kodekeluaran}
def enter_state(self, state):
    self.state = state
    if state == IDLE_STATE:
        self.idle_timer = 0.0
    elif state == PATROL_STATE:
        self.waypoints_visited = 0
    elif state == CHASE_STATE:
        self.lost_timer = 0.0
    self.transitions.append((self.frames, state))
\end{kodekeluaran}

Perhatikan tiga hal yang dikerjakan ketika keadaan dimasuki. Pertama, keadaan
ditetapkan. Kedua, penghitung yang berkaitan dengan keadaan itu dikembalikan ke
nol; tanpa itu, pewaktu dari kunjungan sebelumnya masih tersisa dan perpindahan
berikutnya terjadi terlalu cepat. Ketiga, perpindahan dicatat, dan pencatatan
itulah yang dipakai untuk pengujian pada bagian berikutnya.

\begin{catatabox}
Perhatikan baris terakhir pada fungsi \code{enter_state}. Mencatat perpindahan
keadaan di dalam program bukan sekadar keperluan pengujian; catatan yang sama
juga berguna ketika program dijalankan sungguhan, misalnya untuk menyimpan
riwayat perilaku NPC demi keperluan penelusuran kesalahan. Kebiasaan mencatat
keadaan bukan hal yang asing pada perangkat lunak: ia dipakai pada peladen,
pengendali perangkat, dan mesin permainan.
\end{catatabox}

\subsection{Penghalang pada Batas Jangkauan Serang}

Satu hal yang diwarisi dari Bab~VIII adalah penghalang pada batas. Perhatikan
syarat perpindahan dari keadaan menyerang: pemain harus keluar dari jangkauan
serang \textbf{dikalikan 1,2}, bukan jangkauan serang itu sendiri.

\begin{kodekeluaran}
if self.state == ATTACK_STATE:
    if self.distance > self.attack_range * 1.2:
        return CHASE_STATE
    return ATTACK_STATE
\end{kodekeluaran}

Tanpa penghalang itu, kedua syarat perpindahan pada pasangan keadaan mengejar
dan menyerang akan benar pada saat yang sama ketika jaraknya tepat sama dengan
jangkauan serang. Akibatnya NPC berganti keadaan pada setiap bingkai, dan
gejalanya terlihat sebagai NPC yang berkedip dan tersendat di tepi jangkauan.
Cara yang sama telah dipakai pada Bab~VIII untuk batas persepsi, dan alasan
keperluannya pun sama.
